%% ch09 --- Three equations for the weather: the Lorenz attractor
%%
%% Stub, generated by tools/gen_stubs.py from tools/chapters.json.
%% The brief below is the contract for this chapter: what it must argue, what
%% it must buy the reader, and what it must NOT take from its neighbours.
%% Writing the chapter means deleting the stub block below (the one that opens
%% on the next \chapter line) and replacing it with the chapter text -- after
%% which this file is never regenerated.
%% If the brief turns out to be wrong once the code has been run, change it in
%% the manifest and record the contradiction in CLAUDE.md under *Resolved
%% questions*.

\chapter{Three equations for the weather: the Lorenz attractor}
\label{ch:lorenz}

\chapterstub{%
Argument: chaos is not a quirk of stepping rules; it lives in continuous
time too, and Lorenz found it in 1963 in three equations cut down from a
model of convection (sigma = 10, rho = 28, beta = 8/3), in his paper
Deterministic Nonperiodic Flow. Integrated with Chapter 4's RK4, the state
wanders round two lobes, switching unpredictably between them, and never
repeats \dash{} three numbers is exactly what Chapter 4's theorem said chaos
in a flow needs. Two starts a hair apart diverge, and the rate is a Lyapunov
exponent of about 0.9 per time unit, measured here. And yet the trajectory
never leaves a bounded butterfly-shaped set: the attractor. That is the
chapter's hinge and the book's second big idea: the INDIVIDUAL future is
unpredictable, but WHERE the system lives and HOW OFTEN it visits each
region is perfectly predictable. Statistics of the attractor \dash{} the
fraction of time on each lobe, the average height \dash{} are stable
numbers, the same from almost any start. Payoff: the reader can integrate
the Lorenz system, measure its divergence, and compute a statistic of the
attractor that does not depend on the start. Traps to elicit: a chaotic
system can go anywhere; if the path is unpredictable nothing about it is
predictable; chaos needs many variables. Listings: RK4 on the Lorenz system;
divergence of two runs; lobe statistics from different starts. Labs:
implement the Lorenz derivative; measure a lobe fraction. Figures: the
attractor in projection; two runs diverging; a histogram that does not care
about the start.

\medskip
\textbf{Word budget:} 3000.
}
